\documentclass[11pt]{article}

\usepackage{series}
\usepackage[T1]{fontenc}
\usepackage[utf8]{inputenc}
\usepackage{lmodern}
\usepackage{microtype}
\usepackage{amsmath,amssymb}
\usepackage{booktabs}
\usepackage{enumitem}
\usepackage{longtable}

\newcommand{\cO}{\mathcal{O}}
\newcommand{\cH}{\mathcal{H}}
\newcommand{\TT}{\mathrm{TT}}
\newcommand{\BRST}{\mathrm{BRST}}
\newcommand{\BV}{\mathrm{BV}}

\title{Modal Triplet Theory: Perturbative Coherent-Sector Quantum Gravity and the Heterotic UV-Completion Boundary}
\author{Peter Nero}
\date{Sixth edition\\July 2026}
\hypersetup{
  pdftitle={Modal Triplet Theory: Perturbative Coherent-Sector Quantum Gravity and the Heterotic UV-Completion Boundary},
  pdfsubject={Current low-energy quantum-gravity parity tier and q79 heterotic completion boundary}
}

\begin{document}
\maketitle

\begin{abstract}
This paper separates three levels of the current Modal Triplet Theory (MTT)
quantum-gravity program. First, an exact finite-carrier result supplies
internal transverse-traceless support, while a separate controlled reduction
recovers the four-dimensional Einstein action only under declared
higher-dimensional action, truncation, gap, and normalization hypotheses.
Second, once the resulting local Einstein--matter effective action, Wilson
coefficients, state, gauge fixing, and scheme are supplied, standard
background-field and BV/BRST methods give fixed-order perturbative quantum-GR
equivalence. This is an imported parity construction, not a derivation of the
quantum measure or Wilson coefficients from MTT. The two-derivative theory is
not ultraviolet complete: its counterterm order grows with loop number and
the two-loop pure-gravity divergence is nonzero. Third, the selected
compatible ultraviolet route is the q79 Fu--Yau heterotic branch. A
conditional inheritance theorem gives absence of local ultraviolet
divergences at each fixed genus and multiplicity if an exact modular,
anomaly-free heterotic worldsheet theory with factorization, GSO projection,
and BV vertices is supplied. The current q79 worldsheet contract has five of
twelve rows available, two partial, and five open. The physical nonpullback
holomorphic/HYM bundles, complete Green--Schwarz representative, exact
worldsheet SCFT, string-field realization, infrared completion, all-genus
control, and nonperturbative definition remain open. The resulting claim is
therefore low-energy quantum-gravity parity plus a sharply specified
heterotic completion program, not completed ultraviolet quantum gravity.
\end{abstract}

\section*{Revision note for this edition}
\addcontentsline{toc}{section}{Revision note for this edition}
\begin{description}
\item[Supersedes.] \emph{Modal Triplet Theory: Perturbative Coherent-Sector
Quantum Gravity and the Heterotic UV-Completion Boundary}, version~5.
\item[Reason.] Version~5 correctly discovered the Gaussian
positivity no-go and loop-rank obstruction, but it repeated those results in
several appendices, retained an unselected schedule as a derivation of the
SPT endpoint, and sometimes described imported perturbative QFT machinery as
though it were selected by MTT.
\item[Resolution.] This edition assigns the SPT mathematics to the companion
Euclidean TT paper, states the low-energy result at its exact fixed-order
parity tier, presents all twelve q79 worldsheet gates, and separates
fixed-genus ultraviolet inheritance from all-genus and nonperturbative
completion.
\item[Retained result.] The current chain supports an exact internal TT
carrier, conditional classical GR reduction, fixed-order quantum-GR EFT
parity, selection of the q79 heterotic route as the primary compatible UV
candidate, and a conditional fixed-genus inheritance theorem.
\item[Remaining boundary.] MTT has not yet selected the physical
higher-derivative Wilson values, complete Lorentzian quantum state and
measure, physical nonpullback heterotic bundles, exact q79 worldsheet theory,
or nonperturbative ultraviolet completion.
\end{description}

\tableofcontents

\section{Purpose, Orientation, and Claim Tiers}

In plain language, the present construction has a low-energy bridge and a
candidate high-energy bridge. The low-energy bridge says that if MTT supplies
the same effective gravitational data used by quantum GR, then the standard
quantization machinery yields the same fixed-order observables. The
high-energy bridge says that if the selected q79 geometry can be completed
into an exact heterotic worldsheet theory, it inherits the string
worldsheet's fixed-genus ultraviolet behavior. Neither conditional bridge
manufactures its missing input.

The logical tiers are:
\begin{center}
\begin{tabular}{p{0.31\linewidth}p{0.23\linewidth}p{0.34\linewidth}}
\toprule
Tier & Current status & Meaning\\
\midrule
Finite internal TT carrier & exact on the declared finite branch &
Internal helicity-two support, not yet a normalized Lorentzian graviton\\
Four-dimensional Einstein reduction & conditional & Requires an
higher-dimensional Einstein sector, controlled truncation, and stabilized
non-Einstein modes\\
Interacting quantum-GR EFT & fixed-order parity & Standard QFT structure is
imported after action, Wilson data, state, and scheme are declared\\
SPT Gaussian filter & conditional Euclidean model & Useful graph regulator;
not selected physical propagation\\
q79 heterotic route & selected compatible route & The route is selected, but
its physical worldsheet object is incomplete\\
Fixed-genus string UV inheritance & conditional theorem & Applies only after
the exact worldsheet hypotheses are met\\
All-genus or nonperturbative QG & open & No convergence or constructive
definition has been supplied\\
\bottomrule
\end{tabular}
\end{center}

The word \emph{parity} is important. It means equality to an already defined
target theory after matching all target inputs at a declared order. It is a
valid equivalence statement, but it is not a no-input prediction or a
derivation of those inputs.

\section{The Low-Energy Coherent-Sector Tier}

\subsection{What MTT supplies before quantization}

The finite proto-spinor/GR program proves an exact internal TT support
identity on the selected finite carrier. It identifies the intended
helicity-two internal block and removes an ambiguity that was formerly left
as an ansatz. It does not by itself provide:
\begin{itemize}
\item a Lorentzian principal symbol and physical massless pole;
\item the Newton or Planck normalization;
\item universal conserved coupling to the full stress tensor;
\item the causal state or asymptotic prescription; or
\item the interacting higher-derivative Wilson coefficients.
\end{itemize}

The companion paper \emph{Controlled Coherent Reduction to Four-Dimensional
Einstein Gravity} supplies a separate analytic bridge. On a
\(4+6\)-dimensional geometry, it gives a contraction/Schur--Feshbach estimate
for eliminating heavy modes, provided that the higher-dimensional Einstein
sector, gap, small coupling, truncation, and zero-mode audit are already
available. The shared circle is counted once as common holonomy data and is
not identified with compact physical time.

\subsection{Fixed-order quantum-GR parity}

Let \(S_{\rm EFT}^{(N)}\) denote the local diffeomorphism-invariant
Einstein--matter effective action through a declared loop and derivative
order \(N\). It includes the Einstein term, cosmological term, matter sector,
and the finite list of local higher-curvature operators required at that
order. Let \(\sigma\) denote the causal state, gauge fixing, scale, and
renormalization scheme. The current parity construction is summarized by
\begin{equation}
 \cO_{\rm MTT}^{(N)}
 =
 \operatorname{Readout}\circ
 \operatorname{Green}_{\sigma}\circ
 Q_{\rm GR,EFT}^{(N)}\circ
 E_{q79}.
 \label{eq:composition}
\end{equation}
Here \(E_{q79}\) embeds the declared coherent source data into the effective
field variables, \(Q_{\rm GR,EFT}^{(N)}\) is the standard fixed-order
background-field/BV quantization, and the final maps construct the selected
observables.

\begin{proposition}[Fixed-order parity statement]
\label{prop:parity}
If \(E_{q79}\) supplies the same renormalized action
\(S_{\rm EFT}^{(N)}\), Wilson coefficients, gauge fixing, state, scale, and
scheme as quantum GR EFT through order \(N\), then the renormalized generating
functionals and gauge-invariant observables obtained from
\eqref{eq:composition} agree through order \(N\).
\end{proposition}

\begin{proof}
After the stated data are identified, both constructions apply the same
renormalized functional and the same observable maps. Equality follows by
composition. The proposition proves a typed embedding/equivalence result; it
does not select the matched data.
\end{proof}

\paragraph{What this means.}
The proposition checks that the MTT carrier can host the standard
fixed-order quantum-GR calculation without changing its predictions. Its
scientific content is compatibility and typed transport. Predictive
improvement would require MTT to emit some of the matched Wilson values,
state data, or normalization before the target calculation is imported.

On a globally hyperbolic background with a Hadamard state, local covariant
time-ordered products and perturbative observables can be constructed using
Epstein--Glaser or perturbative algebraic QFT methods
\cite{BrunettiFredenhagen2000,HollandsWald2001,HollandsWald2002,Rejzner2016}.
The BV induction restores the quantum master equation order by order when
the relevant local ghost-number-one anomaly class vanishes
\cite{BV1981,BV1983,BarnichBrandtHenneaux2000}. These are standard imported
theorems. The current MTT source does not yet derive the full measure,
Hadamard state, anomaly cohomology, or fixed-nonzero-coupling completion.

\subsection{Why fixed-order predictivity is not UV completion}

For a connected Einstein graph with \(L\) loops, \(V\) two-derivative
vertices, and \(I\) internal lines, the superficial derivative order is
\[
 D=4L+2V-2I.
\]
Using \(L=I-V+1\) gives
\begin{equation}
 D=2L+2.
 \label{eq:power}
\end{equation}
For example, tree, one-loop, and two-loop graphs require operators of
superficial derivative order two, four, and six respectively. The
Riemann-cubed counterterm at two loops is therefore the first explicit
worked case showing why the two-derivative parameter set cannot remain
closed.
At every fixed loop order only finitely many local counterterms are needed,
which is precisely the effective-field-theory notion of predictivity
\cite{Donoghue1994}. Across all loop orders, however, the derivative order is
unbounded. Pure gravity is on-shell finite at one loop for zero cosmological
constant in the standard sense, but it has a nonzero two-loop
Riemann-cubed divergence \cite{tHooftVeltman1974,GoroffSagnotti1985}.

Consequently, the two parameters of the two-derivative action do not define
an all-scale interacting quantum theory. The current MTT parity theorem
allows a finite set of declared Wilson values at any chosen order; it does
not emit their complete tower.

\section{Why SPT Does Not Close the Ultraviolet Problem}

The companion paper \emph{An SPT-Filtered Euclidean TT Model and Its
Conditional Perturbative Properties} owns the SPT theorems. Its conclusions
can be summarized without repeating their proofs:
\begin{enumerate}
\item an internal spectral gap controls internal inversion but does not
produce an external Gaussian multiplier;
\item a positive proper-time endpoint is an additional filter or schedule
assumption;
\item Gaussian-damped lines control all \(L\) loop directions exactly when
their loop-momentum matrix \(C_\Gamma\) has
\(\operatorname{rank}C_\Gamma=L\), equivalently when unfiltered edges form a
forest; and
\item a nonzero positive Kallen--Lehmann/Stieltjes propagator cannot have
permanent Gaussian decay in the same Euclidean spectral variable.
\end{enumerate}

These results leave SPT useful as a Euclidean regulator or exploratory
effective filter. They withdraw the former route
\[
 \text{internal gap}
 \Longrightarrow
 \text{physical Gaussian propagator}
 \Longrightarrow
 \text{OS positivity, causality, and unitarity}.
\]
Each arrow needs an independent theorem, and the first arrow is false in
general. The exact finite TT eigenvalue recorded by the internal carrier is
dimensionless until a physical units-and-normalization map is supplied; it is
not automatically the Planck cutoff.

\section{The Selected q79 Heterotic Route}

\subsection{Why this is the primary compatible candidate}

The q79 program selects a Fu--Yau-type heterotic route rather than the SPT
filter as its primary compatible ultraviolet candidate. Fu--Yau geometry
provides established existence machinery for non-Kahler heterotic
compactifications with torsion under specific bundle, stability, and anomaly
conditions \cite{FuYau2006}. Torsion linear sigma models provide a related
worldsheet framework \cite{AdamsErnebjergLapan2006}. These external theorems
do not say that every topological MTT packet is automatically a physical
heterotic vacuum.

The current q79 records include:
\begin{itemize}
\item a time-oriented discrete branch label \(q=79\), with conjugate
\(369\) modulo \(448\);
\item a degree-two K3 incidence GLSM base and rank-one Fu--Yau source data;
\item curvature-level Green--Schwarz/Bianchi information;
\item finite gerbe, torsion-phase, and modular-orbit data;
\item a smooth topological visible \(SU(3)\) candidate with
\(c_2=9u\) and \(c_3=\pm6\); and
\item partial twisted-spectral and modular-character calculations.
\end{itemize}

Topological existence and Hodge admissibility do not yet produce the
physical nonpullback holomorphic bundle, a common positive HYM chamber,
connections, or the complete local worldsheet anomaly data. Those are the
physical endpoint obligations, not optional polish.

\subsection{Conditional fixed-genus inheritance}

\begin{theorem}[Heterotic fixed-genus UV inheritance]
\label{thm:inheritance}
Assume the selected q79 background is completed to an exact anomaly-free
modular heterotic \((0,2)\) worldsheet theory with a tachyon-free GSO
projection, correct factorization, and quantum-BV string vertices. Then each
fixed-genus, fixed-multiplicity on-shell amplitude has no local
point-particle ultraviolet divergence. Possible degeneration-boundary
singularities are governed by physical factorization channels and require
the separate tadpole, vacuum-shift, and infrared prescriptions.
\end{theorem}

\begin{proof}[Reason for the inheritance]
Worldsheet moduli are integrated over a modular fundamental domain rather
than over an unrestricted point-particle proper-time region. Short-distance
operator collisions are controlled by the worldsheet OPE and BRST
factorization, while modular equivalence removes duplicate ultraviolet
regions. Degenerations of the Riemann surface represent long tubes and
physical on-shell propagation, so their divergences are infrared or tadpole
questions rather than local ultraviolet counterterms. This is the standard
perturbative string framework applied conditionally to the q79 background
\cite{Polchinski1998v1,SenSuperstringPerturbation,SenUVIR2015}.
\end{proof}

The theorem is coefficientwise in genus and multiplicity. It does not prove
convergence of the genus expansion, existence of a preferred vacuum, or a
nonperturbative definition.

\subsection{The twelve-row worldsheet contract}

The present contract is not a vague list of future work. It has twelve named
rows, with five available, two partial, and five open:

\begin{longtable}{p{0.07\linewidth}p{0.48\linewidth}p{0.32\linewidth}}
\toprule
Row & Required object & Current status\\
\midrule
\endfirsthead
\toprule
Row & Required object & Current status\\
\midrule
\endhead
W1 & Time-oriented q79 target branch & Available\\
W2 & Fu--Yau Mukai charge and Green--Schwarz/Bianchi sector & Available\\
W3 & Visible curvature-level Green--Schwarz cancellation & Available at
curvature tier\\
W4 & Critical heterotic central-charge balance & Available universally\\
W5 & q79 low-energy GR and quantum-EFT limit & Available at declared parity
tier\\
W6 & Global Deligne gerbe and Freed--Witten condition on the full visible
cycle set & Open; finite representative only\\
W7 & All-orders-in-\(\alpha'\) q79 target background & Open; current
background is first-order\\
W8 & Exact q79 heterotic \((0,2)\) SCFT or all beta functions & Partial;
base GLSM, local anomaly, and topological candidate exist, but the physical
bundle and IR SCFT do not\\
W9 & q79 modular-invariant GSO partition function and factorization &
Partial; finite torsion module and seven-seed induction exist, analytic
characters and GSO remain open\\
W10 & q79-specific string-field vertices and BV master action & Open;
standard framework exists but is not instantiated on q79\\
W11 & Tadpole, vacuum-shift, infrared, and soft-state completion & Open\\
W12 & All-genus convergence or a nonperturbative definition & Open\\
\bottomrule
\end{longtable}

Rows W8 and W9 are genuine advances, but partial rows do not count as
available rows. The current authoritative total remains \(5/12\), with two
additional partials. Standard heterotic string-field BV machinery is known
\cite{SenHeteroticBV}; importing its existence is not the same as building
the q79 vertices.

\section{Branch Orientation, Conjugation, and Double Traversal}

The q79 and \(369\) labels form a conjugate pair modulo \(448\):
\[
 79+369\equiv0\pmod{448}.
\]
Current MTT results select the retarded \(q79/F/m1\) representative at the
orientation level. They do not prove uniqueness of the global carrier
measure, derive the initial state, or require the conjugate branch to be
well-behaved.

A Schwinger--Keldysh or closed-time-path description doubles fields along
forward and backward contour legs. Given a normalized density operator and
unitary evolution, equal sources obey \(Z[J,J]=1\)
\cite{HaehlLoganayagamRangamani2017,CrossleyGloriosoLiu2017}. This identity
expresses cancellation on a closed contour; it does not select the arrow of
time or prove that a literal double traversal of the MTT circle returns the
universe to a state with no space, time, or gravity.

The shared MTT circle is therefore treated as common phase or holonomy data.
Physical time is a noncompact Lorentzian ordering variable, possibly related
to a lift of phase data only after an additional theorem.

\section{What Has Been Achieved and What Would Close the Program}

\subsection{Achievements at their declared tiers}

\begin{itemize}
\item The internal TT sector is selected exactly on the finite carrier.
\item A controlled analytic reduction to four-dimensional Einstein gravity
is available under explicit hypotheses.
\item Quantum GR EFT parity is closed at each fixed declared order after
matching the action, Wilson data, state, gauge fixing, scale, and scheme.
\item The finite internal carrier and the old SPT shortcut are excluded as
automatic all-scale ultraviolet regulators.
\item The q79 Fu--Yau heterotic route is selected as the primary compatible
UV route.
\item Fixed-genus UV inheritance is proved conditionally on the exact
worldsheet contract.
\item The worldsheet frontier is counted explicitly as five available, two
partial, and five open rows.
\end{itemize}

\subsection{The shortest honest closure sequence}

The remaining work has a dependency order:
\begin{enumerate}
\item \textbf{Physical Hull--Strominger endpoints.} Construct the visible and
hidden nonpullback holomorphic bundles, common HYM chamber, connections, and
local Bianchi/anomaly representative.
\item \textbf{Complete W6--W9.} Promote the global gerbe, all-order
background, exact \((0,2)\) SCFT, analytic modular characters, GSO
projection, and factorization on that same source.
\item \textbf{Instantiate W10--W11.} Build q79-specific string-field/BV
vertices and the tadpole, vacuum-shift, infrared, and soft-state completion.
\item \textbf{Connect low and high energy.} Derive the four-dimensional
Wilson coefficients, Newton normalization, universal stress response, and
selected state from the same compactification.
\item \textbf{Nonperturbative exit.} Establish all-genus control or another
well-defined nonperturbative completion, with a positive physical observable
sector and comparison map.
\end{enumerate}

The first item controls the rest: without the physical bundles, worldsheet
rows cannot all refer to one selected background. The final item cannot be
inferred merely by adding more fixed-genus calculations.

\section{Conclusion}

The current MTT quantum-gravity result is substantial but tiered. It supplies
an exact internal TT carrier, a conditional Einstein reduction, and
fixed-order quantum-GR EFT parity. It also identifies the q79 Fu--Yau
heterotic construction as the primary compatible ultraviolet route and
proves the correct conditional fixed-genus inheritance statement.

It does not yet supply a UV-complete theory. The former Gaussian shortcut is
a conditional Euclidean model, the low-energy quantum theory imports
standard EFT quantization after matching its inputs, and the q79 worldsheet
contract remains \(5/12\) with two partial rows. The path forward is therefore
not to rename these tiers as closure. It is to construct one physical
visible--hidden Hull--Strominger source, complete its worldsheet and BV
contracts, derive the low-energy normalization and Wilson data from that same
source, and then solve the all-genus or nonperturbative problem.

\bibliographystyle{unsrt}
\bibliography{main}

\end{document}
